\section{Tables and images}\label{sec:tables-images}
\subsection{A basic table}
Table~\ref{tab:tools} uses booktabs rules for a clean academic style.
\begin{table}[H]\centering
\caption{Example tools used in a documentation workflow.}\label{tab:tools}
\begin{tabular}{lll}\toprule
Tool & Main purpose & Source format\\\midrule
LaTeX & Typeset reports & \texttt{.tex}\\
Doxygen & Generate API documentation & C++ comments\\
Beamer & Create slide presentations & \texttt{.tex}\\
Gnuplot & Plot numerical data & \texttt{.gp}\\\bottomrule
\end{tabular}
\end{table}
\subsection{Merged cells and dashed borders}
The \texttt{multirow} and \texttt{multicolumn} commands merge cells. The \texttt{arydshln} package supports dashed rules.
\begin{table}[H]\centering
\caption{A compact table with merged cells and dashed rules.}\label{tab:merged}
\begin{tabular}{|l|c|c|}\hline
\multirow{2}{*}{\textbf{Area}} & \multicolumn{2}{c|}{\textbf{Evidence}}\\
\cdashline{2-3}
& Source file & Output\\\hline
Documentation & C++ comments & HTML pages\\\hline
Presentation & Beamer source & PDF slides\\\hline
\end{tabular}
\end{table}
\subsection{Tables spanning multiple pages}
Use \texttt{longtable} when rows may continue over a page boundary. Its header can repeat on later pages.
\begin{longtable}{p{0.22\linewidth}p{0.66\linewidth}}
\caption{Example long-table content.}\label{tab:long}\\
\toprule\textbf{Stage}&\textbf{Description}\\\midrule
\endfirsthead
\toprule\textbf{Stage}&\textbf{Description (continued)}\\\midrule
\endhead
Planning & Identify reader, purpose, scope, constraints and expected output.\\
Audience analysis & Record what the reader knows and what they need to do.\\
Outline & Arrange sections and decide which diagrams, tables or examples are necessary.\\
Drafting & Organise information and write a first version.\\
Evidence & Check facts, measurements and code examples against reliable sources.\\
Figures & Number diagrams, add captions and refer to them in text.\\
References & Cite external ideas, data and quotations consistently.\\
Review & Check correctness, clarity, consistency and accessibility.\\
Publishing & Export to the required format and preserve editable sources.\\
Maintenance & Update the document when the system or requirements change.\\
\bottomrule
\end{longtable}
\subsection{Adding, scaling and rotating images}
The \texttt{graphicx} package provides \texttt{width}, \texttt{scale} and \texttt{angle} options:
\begin{verbatim}
\includegraphics[width=0.65\linewidth]{diagram.png}
\includegraphics[scale=0.5]{diagram.png}
\includegraphics[width=0.5\linewidth,angle=90]{diagram.png}
\end{verbatim}
Keep proportions unless there is a clear reason to stretch an image. Add a caption and label.
